
\documentclass{article}
%%%%%%%%%%%%%%%%%%%%%%%%%%%%%%%%%%%%%%%%%%%%%%%%%%%%%%%%%%%%%%%%%%%%%%%%%%%%%%%%%%%%%%%%%%%%%%%%%%%%%%%%%%%%%%%%%%%%%%%%%%%%%%%%%%%%%%%%%%%%%%%%%%%%%%%%%%%%%%%%%%%%%%%%%%%%%%%%%%%%%%%%%%%%%%%%%%%%%%%%%%%%%%%%%%%%%%%%%%%%%%%%%%%%%%%%%%%%%%%%%%%%%%%%%%%%
%TCIDATA{OutputFilter=LATEX.DLL}
%TCIDATA{Version=5.50.0.2953}
%TCIDATA{<META NAME="SaveForMode" CONTENT="1">}
%TCIDATA{BibliographyScheme=Manual}
%TCIDATA{Created=Thursday, May 21, 2026 20:22:52}
%TCIDATA{LastRevised=Friday, May 22, 2026 15:27:33}
%TCIDATA{<META NAME="GraphicsSave" CONTENT="32">}
%TCIDATA{<META NAME="DocumentShell" CONTENT="Standard LaTeX\Blank - Standard LaTeX Article">}
%TCIDATA{CSTFile=40 LaTeX article.cst}

\newtheorem{theorem}{Theorem}
\newtheorem{acknowledgement}[theorem]{Acknowledgement}
\newtheorem{algorithm}[theorem]{Algorithm}
\newtheorem{axiom}[theorem]{Axiom}
\newtheorem{case}[theorem]{Case}
\newtheorem{claim}[theorem]{Claim}
\newtheorem{conclusion}[theorem]{Conclusion}
\newtheorem{condition}[theorem]{Condition}
\newtheorem{conjecture}[theorem]{Conjecture}
\newtheorem{corollary}[theorem]{Corollary}
\newtheorem{criterion}[theorem]{Criterion}
\newtheorem{definition}[theorem]{Definition}
\newtheorem{example}[theorem]{Example}
\newtheorem{exercise}[theorem]{Exercise}
\newtheorem{lemma}[theorem]{Lemma}
\newtheorem{notation}[theorem]{Notation}
\newtheorem{problem}[theorem]{Problem}
\newtheorem{proposition}[theorem]{Proposition}
\newtheorem{remark}[theorem]{Remark}
\newtheorem{solution}[theorem]{Solution}
\newtheorem{summary}[theorem]{Summary}
\newenvironment{proof}[1][Proof]{\noindent\textbf{#1.} }{\ \rule{0.5em}{0.5em}}
\input{tcilatex}

\begin{document}


MODELO MATEM\'{A}TICO DEL SISTEMA 


Fase 2: Calcular de forma an\'{a}litica la funci\'{o}n de transferencia, el
error en estado estacionario y el modelo de ecuaciones
integro-diferenciales; determinar la estabilidad en lazo abierto del caso y
del control utilizando los valores num\'{e}ricos para cada sistema;
verificar que las respuestas en lazo abierto obtenida en Simulink de los
modelos matem\'{a}ticos coincidan con la respuesta del circuityo RLC en
SIMSCAPE/Simulink. 

\bigskip 

2.1 ECUACIONES INTEGRO-DIFERENCIALES DEL CIRCUITO 

\begin{eqnarray*}
V_{e}(t) &=&R_{1}i_{1}(t)+\frac{1}{C}\int (i_{1}(t)-i_{2}(t))dt \\
\frac{1}{C}\int (i_{1}(t)-i_{2}(t))dt &=&R_{2}i_{2}(t)+L\frac{di_{2}(t)}{dt}
\\
V_{s}(t) &=&L\frac{di_{2}(t)}{dt}
\end{eqnarray*}

2.2 FUNCI\'{O}N DE TRANSFERENCIA 

Aplicando la Transformada de Laplace: 

\begin{eqnarray*}
V_{e}(s) &=&R_{1}I_{1}(s)+\frac{1}{Cs}[I_{1}(s)-I_{2}(s)] \\
\frac{1}{Cs}[I_{1}(s)-I_{2}(s)] &=&R_{2}I_{2}(s)+LsI_{2}(s) \\
V_{s}(s) &=&LsI_{2}(s)
\end{eqnarray*}

Despejamos I1 de la segunda ecuaci\'{o}n: 

\begin{eqnarray*}
\frac{1}{Cs}[I_{1}(s)-I_{2}(s)] &=&(R_{2}+Ls)I_{2}(s) \\
I_{1}(s)-I_{2}(s) &=&Cs(R_{2}+Ls)I_{2}(s) \\
I_{1}(s) &=&I_{2}(s)+Cs(R_{2}+Ls)I_{2}(s) \\
I_{1}(s) &=&[1+CR_{2}s+CLs^{2}]I_{2}(s)
\end{eqnarray*}

\bigskip 

Entonces: 

\[
V_{e}(s)=R_{1}I_{1}(s)+(R_{2}+Ls)I_{2}(s)
\]

Sustituimos en I1(s): 

\[
V_{e}(s)=R_{1}[1+CR_{2}s+CLs^{2}]I_{2}(s)+(R_{2}+Ls)I_{2}(s)
\]

Factorizamos I2(s):

\[
V_{e}(s)=[R_{1}+R_{1}CR_{2}s+R_{1}CLs^{2}+R_{2}+Ls]I_{2}(s)
\]

Ordenamos: 

\[
V_{e}(s)=[R_{1}CLs^{2}+(L+R_{1}R_{2}C)s+(R_{1}+R_{2})]I_{2}(s)
\]

Despejamos I2(s) de la ecu. 3: 

\[
I_{2}(s)=LsI2(s)
\]

Por lo tanto, queda: 

\[
I_{2}(s)=\frac{V_{s}(s)}{Ls}
\]

Sustituimos I2 de la ecuaci\'{o}n de entrada: 

\[
V_{e}(s)=[R_{1}CLs^{2}+(L+R_{1}R_{2}C)s+(R_{1}+R_{2})]\frac{V_{s}(s)}{Ls}
\]

Funci\'{o}n de transferencia: 

\[
G(s)=\frac{V_{s}(s)}{V_{e}(s)}=\frac{Ls}{%
R_{1}CLs^{2}+(L+R_{1}R_{2}C)s+(R_{1}+R_{2})}
\]

\bigskip 

2.3 ERROR EN ESTADO ESTACIONARIO 


El error en estado estacionario se define como la diferencia entre la
entrada y la salida de un sistema cuando el l\'{\i}mite en el tiempo tiende
a infinito, este an\'{a}lisis solo es \'{u}til para sistemas estables, por
este motivo se debe determinar antes la estabilidad del sistema; el error en
estado estacionario para sistemas en lazo abierto este dado por: 

\[
E(s)=R(s)[1-G(s)]
\]

El error en estado estacionario en lazo abierto se calcula como: 

\[
e_{ss}=\lim_{s\ \rightarrow 0}\left[ \frac{1}{s}(1-G(s))\right] 
\]

Entonces: 

\[
e_{ss}=\lim_{s\ \rightarrow 0}[1-G(s)]
\]

2.3.1 Error en estado estacionario para el sistema de control

Valores:

\bigskip\ Control piel sana: R1=10 ohm, R2=6.8 ohm, C=0.1F, L=1H


Sustituimos los valores en la funci\'{o}n de transferencia: 

\begin{eqnarray*}
\frac{V_{s}(s)}{V_{e}(s)} &=&\frac{1s}{%
(0.1)(1)(10)s^{2}+(1+(0.1)(10)(6.8))s+(10+6.8)} \\
\frac{V_{s}(s)}{V_{e}(s)} &=&\frac{s}{s^{2}+7.8s+16.8}
\end{eqnarray*}

Error en estado estacionario de control: 

\begin{eqnarray*}
e_{ss} &=&\lim_{s\ \rightarrow 0}\left( 1-\frac{s}{s^{2}+7.8s+16.8}\right) 
\\
e_{ss} &=&[1-0] \\
e_{ss} &=&1
\end{eqnarray*}

2.3.2 Error en estado estacionario para el sistema del caso 

Valores:

Caso quemadura de segundo grado: R1=3.3 ohm, R2=12 ohm, C=0.068F, L=1.5H

\bigskip 

Sustituimos los valores en la funci\'{o}n de transferencia: 

\bigskip 
\begin{eqnarray*}
G(s) &=&\frac{1.5s}{(0.068)(1.5)(3.3)s^{2}+(1.5+(0.068)(3.3)(12))s+(3.3+12)}
\\
G(s) &=&\frac{1.5s}{0.3366s^{2}+4.1928s+15.3}
\end{eqnarray*}

\bigskip 

Error en estado estacionario del caso: 

\begin{eqnarray*}
e_{ss} &=&\lim_{s\ \rightarrow 0}\left( 1-\frac{1.5}{0.3366s^{2}+4.1928s+15.3%
}\right)  \\
e_{ss} &=&[1-0] \\
e_{ss} &=&1
\end{eqnarray*}

\bigskip 

2. 4 ESTABILIDAD DEL SISTEMA 

Un sistema din\'{a}mico estable es aquel que presenta una respuesta acotada
a una entrada o \ \ \ \ \ pertur baci\'{o}n acotada, por lo tanto, con base
en su respuesta, se clasifican de la siguiente manera: 


Estable con respuesta sobreamortiguada, subamortiguada o cr\'{\i}ticamente
amortiguada. 

Marginalmente estable con respuesta oscilatoria o no amortiguada. 

Inestable con respuesta divergente (la soluci\'{o}n tiende al infinito). 


La estabilidad de un sistema se determina al calcular los polos, es decir,
las ra\'{\i}ces del \ \ \ denominador de su funci\'{o}n de transferencia. 

\bigskip 

2.4.1 Estabilidad del sistema en lazo abierto 


Para analizar la estabilidad en lazo abierto se utiliza el denominador de la
funci\'{o}n de transferencia: 

\[
G(s)=\frac{Ls}{R_{1}CLs^{2}+(L+R_{1}R_{2}C)s+(R_{1}+R_{2})}
\]

\bigskip 

Para determinar la estabilidad se calculan las ra\'{\i}ces del denominador
usando la f\'{o}rmula general: 

\[
s=\frac{-b\pm \sqrt{b^{2}-4ac}}{2a}
\]

Donde:

\begin{eqnarray*}
a &=&R_{1}CL \\
b &=&L+R_{1}R_{2}C \\
c &=&R_{1}+R_{2}
\end{eqnarray*}

\bigskip 

2.4.2 Estabilidad del sistema en lazo abierto para control 

Valores:

Control piel sana: R1=10 ohm, R2=6.8 ohm, C=0.1F, L=1H

\bigskip 

Por lo tanto:

\[
s^{2}+7.8s+16.8=0
\]

Sustituyendo: 

\[
s=\frac{-7.8\pm \sqrt{(7.8)^{2}-4(1)(16.8)}}{2(1)}
\]

Resultados:

\begin{eqnarray*}
s_{1} &=&-3.9+1.2609j \\
s_{2} &=&-3.9-1.2609j
\end{eqnarray*}

\bigskip 

Por lo tanto, como las ra\'{\i}ces \U{e315}i del denominador de la funcion
de trsferencia son complejas conjugadas\ \ \U{e31d}\pm j\U{e323}\ con parte
real negativa \ \ \text{Re}\U{e315}_{i}\TEXTsymbol{<}0, entonces el sistema
control es estable en lazo abierto con una respuesta subamortiguafa. 

\bigskip 

2.4.3 Estabilidad del sistema en lazo abierto para caso 

Valores:

Caso quemadura de segundo grado: R1=3.3 ohm, R2=12 ohm, C=0.068F, L=1.5H

Por lo tanto:

\[
0.3366s^{2}+4.1928s+15.3=0
\]

Sustituyendo: 

\[
s=\frac{-4.1928\pm \sqrt{(4.1928)^{2}-4(0.3366)(15.3)}}{2(0.3366)}
\]

Resultados:

\begin{eqnarray*}
s_{1} &=&-6.2281+2.5815j \\
s_{2} &=&-6.2281-2.5815j
\end{eqnarray*}

Por lo tanto, como las ra\'{\i}ces (\U{3bb}i) del denominador de la funci%
\'{o}n de transferencia son complejas conjugadas (\U{3c3}\pm j\U{3c9}) con
parte real negativa (Re\U{3bb}i\TEXTsymbol{<}0), entonces el sistema del
caso es estable en lazo abierto con una respuesta subamortiguada. 

\bigskip 

CALCULO PARA MODELO EID:

Circuito del sistema 

El circuito del sistema tegumentario est\'{a} formado por: 

\textbullet\ Fuente de entrada: Ve(t) 

\textbullet\ Resistor superficial: R1 

\textbullet\ Capacitor t\'{e}rmico: C 

\textbullet\ Resistor profundo: R2 

\textbullet\ Inductor vascular: L

\bigskip 

Se definen dos corrientes de malla: 

i1(t)=corriente de la malla izquierda 

i2(t)=corriente de la malla derecha

\bigskip 

La corriente que atraviesa el capacitor es:

\[
iC(t)=i1(t)-i2(t)
\]

Malla Izquierda: 

Aplicando LVK: 

\[
V_{e}(t)-V_{R1}(t)-V_{C}(t)=0
\]

Sustituyendo:

\[
V_{e}(t)-R_{1}i_{1}(t)-V_{C}(t)=0
\]

Despejando i1(t):

\begin{eqnarray*}
V_{e}(t)-V_{C}(t) &=&R_{1}i_{1}(t) \\
i_{1}(t) &=&\frac{V_{e}(t)-V_{C}(t)}{R_{1}}
\end{eqnarray*}

Ecuaci\'{o}n del capacitor: 

La relaci\'{o}n corriente-voltaje del capacitor es: 

\bigskip 
\[
i_{C}(t)=C\frac{dV_{C}(t)}{dt}
\]

Pero como Tambien:

\[
i_{1}(t)-i_{2}(t)=C\frac{dV_{C}(t)}{dt}
\]

Se despeja:

\begin{eqnarray*}
\frac{dV_{C}(t)}{dt} &=&\frac{i_{1}(t)-i_{2}(t)}{C} \\
V_{C}(t) &=&\frac{1}{C}\int [i_{1}(t)-i_{2}(t)]dt
\end{eqnarray*}

Malla Derecha

Aplicamos LVK:

\[
V_{C}(t)-V_{R2}(t)-V_{L}(t)=0
\]

El voltaje en el resistor R2 es: 

\[
V_{R2}(t)=R_{2}i_{2}(t)
\]

El voltaje en el inductor es: 

\bigskip 
\[
V_{L}(t)=L\frac{di_{2}(t)}{dt}
\]

Se sustituye:

\bigskip 
\[
V_{C}(t)-R_{2}i_{2}(t)-L\frac{di_{2}(t)}{dt}=0
\]

Se despeja:

\bigskip 
\[
L\frac{di_{2}(t)}{dt}=V_{C}(t)-R_{2}i_{2}(t)
\]

Dividiendo entre L: 

\begin{eqnarray*}
\frac{di_{2}(t)}{dt} &=&\frac{V_{C}(t)-R_{2}i_{2}(t)}{L} \\
i_{2}(t) &=&\int \left( \frac{V_{C}(t)-R_{2}i_{2}(t)}{L}\right) dt
\end{eqnarray*}

\bigskip ECUACION DE SALIDA:

La salida del sistema se toma en el inductor:

\[
V_{s}(t)=V_{L}(t)
\]

\bigskip Y sabiendo lo siguente:

\[
V_{L}(t)=L\frac{di_{2}(t)}{dt}
\]

Y usando una de las ecuaciones anteriores:

\begin{eqnarray*}
L\frac{di_{2}(t)}{dt} &=&V_{C}(t)-R_{2}i_{2}(t) \\
V_{s}(t) &=&V_{C}(t)-R_{2}i_{2}(t)
\end{eqnarray*}

MODELO EID FINAL:

\begin{eqnarray*}
i_{1}(t) &=&\frac{V_{e}(t)-V_{C}(t)}{R_{1}} \\
V_{C}(t) &=&\frac{1}{C}\int [i_{1}(t)-i_{2}(t)]dt \\
i_{2}(t) &=&\int \left( \frac{V_{C}(t)-R_{2}i_{2}(t)}{L}\right) dt \\
V_{s}(t) &=&V_{C}(t)-R_{2}i_{2}(t)
\end{eqnarray*}

\bigskip 

\bigskip 

\bigskip 

\bigskip 

\bigskip 

\end{document}
